% Réplica del CV original (Canva): columna lateral gris con cabecera azul
% marino y columna principal con barras de sección.
%   xelatex "\def\cvlang{es}% Réplica del CV original (Canva): columna lateral gris con cabecera azul
% marino y columna principal con barras de sección.
%   xelatex "\def\cvlang{es}% Réplica del CV original (Canva): columna lateral gris con cabecera azul
% marino y columna principal con barras de sección.
%   xelatex "\def\cvlang{es}% Réplica del CV original (Canva): columna lateral gris con cabecera azul
% marino y columna principal con barras de sección.
%   xelatex "\def\cvlang{es}\input{cv-replica}"   (o: node build.mjs)
\documentclass[10pt]{article}
\usepackage[a4paper,left=77.7mm,right=7mm,top=0mm,bottom=0mm]{geometry}
\input{comun}
\usepackage{fontspec}
\usepackage{xcolor}

\providecommand\cvlang{es}

\defaultfontfeatures{Ligatures=TeX}
\setmainfont{Lora}[
  Path=../assets/fonts/, Extension=.ttf,
  UprightFont=*-Regular, BoldFont=*-Bold,
  ItalicFont=*-Italic, BoldItalicFont=*-BoldItalic]
\newfontfamily\lato{Lato}[
  Extension=.ttf, UprightFont=*-Regular, BoldFont=*-Bold]
\newfontfamily\latoheavy{Lato-Heavy.ttf}
\newfontfamily\latoblack{Lato-Black.ttf}

\definecolor{navy}{HTML}{2C3E50}
\definecolor{lateral}{HTML}{D9D9D9}
\definecolor{chip}{HTML}{5B7D95}
\definecolor{tinta}{HTML}{2B2B2B}

\pagestyle{empty}
\setlength\parindent{0pt}
\hyphenpenalty=10000\exhyphenpenalty=10000
\color{tinta}

\input{contenido-\cvlang}

% --- Main column pieces ----------------------------------------------------
\newcommand\cvbarra[1]{%
  \par\vspace{3.2mm}%
  \noindent\hspace*{-4.7mm}%
  {\setlength\fboxsep{0pt}\colorbox{navy}{\parbox[c][9mm]{140mm}{%
    \hspace{8.5mm}{\color{white}\latoblack\fontsize{12.5}{14}\selectfont
    \raisebox{0.25ex}{\tikz\fill[white] circle (1.05mm);}\hspace{2mm}%
    \addfontfeatures{LetterSpace=10}\MakeUppercase{#1}}}}}%
  \par\vspace{2.6mm}}

\newcommand\cvcentro[1]{%
  \par\vspace{1.2mm}\hspace*{8mm}%
  {\bfseries\fontsize{11.6}{14}\selectfont\addfontfeatures{LetterSpace=5}\MakeUppercase{#1}}\par\vspace{0.6mm}}

\newenvironment{cvlista}
  {\begin{itemize}[leftmargin=13.5mm,labelsep=2.2mm,itemsep=0.45mm,
      label={\normalfont\textbullet}]\raggedright\itshape\fontsize{10.6}{14.6}\selectfont\addfontfeatures{LetterSpace=2.5}}
  {\end{itemize}}

\newcommand\cvestudio[2]{#1 (#2)}

\newcommand\cvpuestotrabajo[3]{%
  \cvcentro{#1}%
  \hspace*{8mm}{\bfseries\itshape\fontsize{10.6}{14.6}\selectfont\addfontfeatures{LetterSpace=1.5}#2}\par
  {\raggedright\leftskip=8mm\itshape\fontsize{10.6}{14.6}\selectfont\addfontfeatures{LetterSpace=2.5}#3\par}%
  \vspace{2.4mm}}

\newcommand\cvdestacado[2]{{\upshape\bfseries #1} #2}
\newcommand\cvrol[2]{{\upshape\bfseries #1} {\upshape #2}}

% --- Sidebar pieces --------------------------------------------------------
\newcommand\cvidioma[3]{}
\newcommand\cvchip[1]{%
  \tikz[baseline=-0.6ex]\node[ellipse,fill=chip,text=white,inner sep=0.3mm,
    minimum height=7.4mm,minimum width=10.4mm,font=\latoblack,
    align=center]{\cvchiplabel{#1}};}
\newcommand\cvchiplabel[1]{%
  \StrLen{#1}[\cvlen]%
  \ifnum\cvlen<4 \fontsize{7.6}{8}\selectfont\MakeUppercase{#1}%
  \else\fontsize{4.4}{4.8}\selectfont
    \ifnum\cvlen>9 \parbox{8.6mm}{\centering\MakeUppercase{#1}}\else\MakeUppercase{#1}\fi\fi}
\usepackage{xstring}
\usetikzlibrary{shapes.geometric}
\newcommand\cvgrupotec[2]{%
  {\latoblack\fontsize{7.3}{9}\selectfont\addfontfeatures{LetterSpace=6}\MakeUppercase{#1}}\par
  \vspace{1.4mm}%
  {\baselineskip=0pt\lineskip=1.5mm\lineskiplimit=0pt\raggedright
   \forcsvlist{\cvchipsep}{#2}\par}\vspace{3.4mm}}
\newcommand\cvchipsep[1]{\cvchip{#1}\hspace{1.3mm}\allowbreak}

\newcommand\fila[3]{% y (mm from top), icon, text
  \node[anchor=base west,inner sep=0] at (8mm,-#1mm) {\color{black}\fontsize{9}{10}\selectfont#2};
  \node[anchor=base west,inner sep=0] at (12.6mm,-#1mm) {\lato\fontsize{8.6}{10}\selectfont\color{black!85}#3};}
\newcommand\enlace[2]{\href{#1}{\underline{#2}}}
\newcommand\titulolateral[2]{%
  \node[anchor=base west,inner sep=0] at (9mm,-#1mm)
    {\latoblack\fontsize{13}{15}\selectfont\addfontfeatures{LetterSpace=10}\color{black!90}\MakeUppercase{#2}};}
\newcommand\separador[1]{\draw[black,line width=0.55pt] (0,-#1mm) -- (72.9mm,-#1mm);}

\begin{document}
\vspace*{7.4mm}\noindent
\begin{tikzpicture}[remember picture,overlay,x=1mm,y=1mm]
  \begin{scope}[shift={(current page.north west)}]
    % Backgrounds.
    \fill[lateral] (0,0) rectangle (72.9,-297.1);
    \fill[navy] (0,0) rectangle (72.9,-77.2);
    \fill[navy] (72.9,0) rectangle (210.1,-52.2);

    % Portrait.
    \begin{scope}
      \clip (36.3,-27.2) circle (21.4);
      \node[inner sep=0] at (36.3,-29.2) {\includegraphics[width=43.6mm]{foto.jpg}};
    \end{scope}
    \draw[white!70!navy,line width=0.7pt] (36.3,-27.2) circle (21.4);

    % Name.
    \node[anchor=base,inner sep=0] at (36.3,-57.3)
      {\latoheavy\fontsize{19.5}{22}\selectfont\addfontfeatures{LetterSpace=4}\color{white}\MakeUppercase{\cvNombre}};
    \node[anchor=base,inner sep=0] at (36.3,-63.4)
      {\latoheavy\fontsize{17}{20}\selectfont\addfontfeatures{LetterSpace=4}\color{white}\MakeUppercase{\cvApellidos}};
    \node[anchor=base,inner sep=0] at (36.3,-71.7)
      {\lato\fontsize{12.4}{14}\selectfont\addfontfeatures{LetterSpace=2}\color{white}\cvPuesto};

    % Personal details.
    \fila{85.7}{\faMapMarker*}{\cvUbicacion}
    \siNacimiento{\fila{95.2}{\faCalendar[regular]}{\cvNacimiento}}
    \separador{101.0}
    \titulolateral{108.0}{\cvt{contacto}}
    \siTelefono{\fila{115.0}{\faPhone*}{\cvTelefono}}
    \fila{124.6}{\faEnvelope[regular]}{\enlace{mailto:\cvEmail}{\cvEmail}}
    \fila{134.2}{\faLinkedin}{\enlace{https://www.linkedin.com/in/\cvLinkedin}{\cvLinkedin}}
    \fila{143.8}{\faGithub}{\enlace{https://github.com/\cvGithub}{\cvGithub}}
    \ifdefined\cvOrcid\fila{153.4}{\aiOrcid}{\enlace{https://orcid.org/\cvOrcid}{\cvOrcid}}\fi
    \separador{160.7}
    \titulolateral{168.6}{\cvt{idiomas}}
    \renewcommand\cvidioma[3]{%
      \node[anchor=west,inner sep=0] at (\cvx,-178.2) {\bandera{#1}};
      \node[anchor=west,inner sep=0,align=left] at (\cvx+10.3,-178.2)
        {\lato\fontsize{8.6}{9.6}\selectfont#2\\#3};
      \pgfmathsetmacro\cvx{\cvx+24.6}\xdef\cvx{\cvx}}
    \def\cvx{6.8}
    \cvIdiomas
    \separador{185.9}
    \titulolateral{195.4}{\cvt{tecnologias}}
    \node[anchor=north west,inner sep=0,text width=66mm] at (5.2,-198.8) {\cvTecnologias};
  \end{scope}
\end{tikzpicture}%
{\color{white}\tolerance=9999\emergencystretch=4em\fontsize{12.4}{18.6}\selectfont\addfontfeatures{LetterSpace=2.5}\cvResumen\par}
\vspace{-1.6mm}
\cvbarra{\cvt{formacion}}
\cvFormacion
\cvbarra{\cvt{aptitudes}}
\cvAptitudes
\cvbarra{\cvt{experiencia}}
\cvExperiencia
\vspace{-2.4mm}
\cvbarra{\cvt{proyectos}}
\cvProyectos
\cvbarra{\cvt{complementaria}}
\cvComplementaria
\end{document}
"   (o: node build.mjs)
\documentclass[10pt]{article}
\usepackage[a4paper,left=77.7mm,right=7mm,top=0mm,bottom=0mm]{geometry}
% Shared by both templates: storage for the content fields, the private data,
% flags and icons. Compile with XeLaTeX (see build.mjs).
\usepackage{etoolbox}
\usepackage{pgfkeys}
\usepackage{graphicx}
\usepackage{tikz}
\usepackage{fontawesome5}
\usepackage{academicons}
\usepackage{enumitem}
\usepackage[hidelinks]{hyperref}

% --- Content fields --------------------------------------------------------
\newcommand\cvnombre[2]{\def\cvNombre{#1}\def\cvApellidos{#2}}
\newcommand\cvpuesto[1]{\def\cvPuesto{#1}}
\newcommand\cvenfoque[1]{\def\cvEnfoque{#1}}
\newcommand\cvresumen[1]{\def\cvResumen{#1}}
\newcommand\cvubicacion[1]{\def\cvUbicacion{#1}}
\newcommand\cvemail[1]{\def\cvEmail{#1}}
\newcommand\cvlinkedin[1]{\def\cvLinkedin{#1}}
\newcommand\cvgithub[1]{\def\cvGithub{#1}}
\newcommand\cvorcid[1]{\def\cvOrcid{#1}}
\newcommand\cvweb[1]{\def\cvWeb{#1}}
\newcommand\cvtelefono[1]{\def\cvTelefono{#1}}
\newcommand\cvnacimiento[1]{\def\cvNacimiento{#1}}

% Section titles: \cvtitulos{clave=Título, ...} and \cvt{clave}.
\pgfkeys{/cvt/.unknown/.code={\csdef{cvt@\pgfkeyscurrentname}{#1}}}
\newcommand\cvtitulos[1]{\pgfqkeys{/cvt}{#1}}
\newcommand\cvt[1]{\csuse{cvt@#1}}

% --- Private data ----------------------------------------------------------
% privado.tex (git-ignored) holds the phone number and birth date. The public
% build (\cvpublico) leaves them out even when the file exists.
\providetoggle{cvpublico}
\ifdefined\cvpublic\toggletrue{cvpublico}\fi
\IfFileExists{privado.tex}{\input{privado.tex}}{}
\iftoggle{cvpublico}{\let\cvTelefono\relax\let\cvNacimiento\relax}{}
\newcommand\siTelefono[1]{\ifdefined\cvTelefono\ifx\cvTelefono\relax\else#1\fi\fi}
\newcommand\siNacimiento[1]{\ifdefined\cvNacimiento\ifx\cvNacimiento\relax\else#1\fi\fi}

% --- Flags -----------------------------------------------------------------
\newcommand\bandera[1]{%
  \begin{tikzpicture}[x=1mm,y=1mm,baseline=0.6mm]
    \ifstrequal{#1}{es}{%
      \fill[red!80!black] (0,0) rectangle (9,6);
      \fill[yellow!85!orange] (0,1.5) rectangle (9,4.5);
      \fill[red!70!black] (2,2.4) rectangle (3.2,3.6);
    }{%
      \fill[blue!50!black] (0,0) rectangle (9,6);
      \begin{scope}
        \clip (0,0) rectangle (9,6);
        \draw[white,line width=1.2mm] (0,0) -- (9,6) (0,6) -- (9,0);
        \draw[red!80!black,line width=0.4mm] (0,0) -- (9,6) (0,6) -- (9,0);
      \end{scope}
      \fill[white] (3.6,0) rectangle (5.4,6) (0,2.1) rectangle (9,3.9);
      \fill[red!80!black] (3.95,0) rectangle (5.05,6) (0,2.45) rectangle (9,3.55);
    }
  \end{tikzpicture}}

\setlist{nosep}

\usepackage{fontspec}
\usepackage{xcolor}

\providecommand\cvlang{es}

\defaultfontfeatures{Ligatures=TeX}
\setmainfont{Lora}[
  Path=../assets/fonts/, Extension=.ttf,
  UprightFont=*-Regular, BoldFont=*-Bold,
  ItalicFont=*-Italic, BoldItalicFont=*-BoldItalic]
\newfontfamily\lato{Lato}[
  Extension=.ttf, UprightFont=*-Regular, BoldFont=*-Bold]
\newfontfamily\latoheavy{Lato-Heavy.ttf}
\newfontfamily\latoblack{Lato-Black.ttf}

\definecolor{navy}{HTML}{2C3E50}
\definecolor{lateral}{HTML}{D9D9D9}
\definecolor{chip}{HTML}{5B7D95}
\definecolor{tinta}{HTML}{2B2B2B}

\pagestyle{empty}
\setlength\parindent{0pt}
\hyphenpenalty=10000\exhyphenpenalty=10000
\color{tinta}

\input{contenido-\cvlang}

% --- Main column pieces ----------------------------------------------------
\newcommand\cvbarra[1]{%
  \par\vspace{3.2mm}%
  \noindent\hspace*{-4.7mm}%
  {\setlength\fboxsep{0pt}\colorbox{navy}{\parbox[c][9mm]{140mm}{%
    \hspace{8.5mm}{\color{white}\latoblack\fontsize{12.5}{14}\selectfont
    \raisebox{0.25ex}{\tikz\fill[white] circle (1.05mm);}\hspace{2mm}%
    \addfontfeatures{LetterSpace=10}\MakeUppercase{#1}}}}}%
  \par\vspace{2.6mm}}

\newcommand\cvcentro[1]{%
  \par\vspace{1.2mm}\hspace*{8mm}%
  {\bfseries\fontsize{11.6}{14}\selectfont\addfontfeatures{LetterSpace=5}\MakeUppercase{#1}}\par\vspace{0.6mm}}

\newenvironment{cvlista}
  {\begin{itemize}[leftmargin=13.5mm,labelsep=2.2mm,itemsep=0.45mm,
      label={\normalfont\textbullet}]\raggedright\itshape\fontsize{10.6}{14.6}\selectfont\addfontfeatures{LetterSpace=2.5}}
  {\end{itemize}}

\newcommand\cvestudio[2]{#1 (#2)}

\newcommand\cvpuestotrabajo[3]{%
  \cvcentro{#1}%
  \hspace*{8mm}{\bfseries\itshape\fontsize{10.6}{14.6}\selectfont\addfontfeatures{LetterSpace=1.5}#2}\par
  {\raggedright\leftskip=8mm\itshape\fontsize{10.6}{14.6}\selectfont\addfontfeatures{LetterSpace=2.5}#3\par}%
  \vspace{2.4mm}}

\newcommand\cvdestacado[2]{{\upshape\bfseries #1} #2}
\newcommand\cvrol[2]{{\upshape\bfseries #1} {\upshape #2}}

% --- Sidebar pieces --------------------------------------------------------
\newcommand\cvidioma[3]{}
\newcommand\cvchip[1]{%
  \tikz[baseline=-0.6ex]\node[ellipse,fill=chip,text=white,inner sep=0.3mm,
    minimum height=7.4mm,minimum width=10.4mm,font=\latoblack,
    align=center]{\cvchiplabel{#1}};}
\newcommand\cvchiplabel[1]{%
  \StrLen{#1}[\cvlen]%
  \ifnum\cvlen<4 \fontsize{7.6}{8}\selectfont\MakeUppercase{#1}%
  \else\fontsize{4.4}{4.8}\selectfont
    \ifnum\cvlen>9 \parbox{8.6mm}{\centering\MakeUppercase{#1}}\else\MakeUppercase{#1}\fi\fi}
\usepackage{xstring}
\usetikzlibrary{shapes.geometric}
\newcommand\cvgrupotec[2]{%
  {\latoblack\fontsize{7.3}{9}\selectfont\addfontfeatures{LetterSpace=6}\MakeUppercase{#1}}\par
  \vspace{1.4mm}%
  {\baselineskip=0pt\lineskip=1.5mm\lineskiplimit=0pt\raggedright
   \forcsvlist{\cvchipsep}{#2}\par}\vspace{3.4mm}}
\newcommand\cvchipsep[1]{\cvchip{#1}\hspace{1.3mm}\allowbreak}

\newcommand\fila[3]{% y (mm from top), icon, text
  \node[anchor=base west,inner sep=0] at (8mm,-#1mm) {\color{black}\fontsize{9}{10}\selectfont#2};
  \node[anchor=base west,inner sep=0] at (12.6mm,-#1mm) {\lato\fontsize{8.6}{10}\selectfont\color{black!85}#3};}
\newcommand\enlace[2]{\href{#1}{\underline{#2}}}
\newcommand\titulolateral[2]{%
  \node[anchor=base west,inner sep=0] at (9mm,-#1mm)
    {\latoblack\fontsize{13}{15}\selectfont\addfontfeatures{LetterSpace=10}\color{black!90}\MakeUppercase{#2}};}
\newcommand\separador[1]{\draw[black,line width=0.55pt] (0,-#1mm) -- (72.9mm,-#1mm);}

\begin{document}
\vspace*{7.4mm}\noindent
\begin{tikzpicture}[remember picture,overlay,x=1mm,y=1mm]
  \begin{scope}[shift={(current page.north west)}]
    % Backgrounds.
    \fill[lateral] (0,0) rectangle (72.9,-297.1);
    \fill[navy] (0,0) rectangle (72.9,-77.2);
    \fill[navy] (72.9,0) rectangle (210.1,-52.2);

    % Portrait.
    \begin{scope}
      \clip (36.3,-27.2) circle (21.4);
      \node[inner sep=0] at (36.3,-29.2) {\includegraphics[width=43.6mm]{foto.jpg}};
    \end{scope}
    \draw[white!70!navy,line width=0.7pt] (36.3,-27.2) circle (21.4);

    % Name.
    \node[anchor=base,inner sep=0] at (36.3,-57.3)
      {\latoheavy\fontsize{19.5}{22}\selectfont\addfontfeatures{LetterSpace=4}\color{white}\MakeUppercase{\cvNombre}};
    \node[anchor=base,inner sep=0] at (36.3,-63.4)
      {\latoheavy\fontsize{17}{20}\selectfont\addfontfeatures{LetterSpace=4}\color{white}\MakeUppercase{\cvApellidos}};
    \node[anchor=base,inner sep=0] at (36.3,-71.7)
      {\lato\fontsize{12.4}{14}\selectfont\addfontfeatures{LetterSpace=2}\color{white}\cvPuesto};

    % Personal details.
    \fila{85.7}{\faMapMarker*}{\cvUbicacion}
    \siNacimiento{\fila{95.2}{\faCalendar[regular]}{\cvNacimiento}}
    \separador{101.0}
    \titulolateral{108.0}{\cvt{contacto}}
    \siTelefono{\fila{115.0}{\faPhone*}{\cvTelefono}}
    \fila{124.6}{\faEnvelope[regular]}{\enlace{mailto:\cvEmail}{\cvEmail}}
    \fila{134.2}{\faLinkedin}{\enlace{https://www.linkedin.com/in/\cvLinkedin}{\cvLinkedin}}
    \fila{143.8}{\faGithub}{\enlace{https://github.com/\cvGithub}{\cvGithub}}
    \ifdefined\cvOrcid\fila{153.4}{\aiOrcid}{\enlace{https://orcid.org/\cvOrcid}{\cvOrcid}}\fi
    \separador{160.7}
    \titulolateral{168.6}{\cvt{idiomas}}
    \renewcommand\cvidioma[3]{%
      \node[anchor=west,inner sep=0] at (\cvx,-178.2) {\bandera{#1}};
      \node[anchor=west,inner sep=0,align=left] at (\cvx+10.3,-178.2)
        {\lato\fontsize{8.6}{9.6}\selectfont#2\\#3};
      \pgfmathsetmacro\cvx{\cvx+24.6}\xdef\cvx{\cvx}}
    \def\cvx{6.8}
    \cvIdiomas
    \separador{185.9}
    \titulolateral{195.4}{\cvt{tecnologias}}
    \node[anchor=north west,inner sep=0,text width=66mm] at (5.2,-198.8) {\cvTecnologias};
  \end{scope}
\end{tikzpicture}%
{\color{white}\tolerance=9999\emergencystretch=4em\fontsize{12.4}{18.6}\selectfont\addfontfeatures{LetterSpace=2.5}\cvResumen\par}
\vspace{-1.6mm}
\cvbarra{\cvt{formacion}}
\cvFormacion
\cvbarra{\cvt{aptitudes}}
\cvAptitudes
\cvbarra{\cvt{experiencia}}
\cvExperiencia
\vspace{-2.4mm}
\cvbarra{\cvt{proyectos}}
\cvProyectos
\cvbarra{\cvt{complementaria}}
\cvComplementaria
\end{document}
"   (o: node build.mjs)
\documentclass[10pt]{article}
\usepackage[a4paper,left=77.7mm,right=7mm,top=0mm,bottom=0mm]{geometry}
% Shared by both templates: storage for the content fields, the private data,
% flags and icons. Compile with XeLaTeX (see build.mjs).
\usepackage{etoolbox}
\usepackage{pgfkeys}
\usepackage{graphicx}
\usepackage{tikz}
\usepackage{fontawesome5}
\usepackage{academicons}
\usepackage{enumitem}
\usepackage[hidelinks]{hyperref}

% --- Content fields --------------------------------------------------------
\newcommand\cvnombre[2]{\def\cvNombre{#1}\def\cvApellidos{#2}}
\newcommand\cvpuesto[1]{\def\cvPuesto{#1}}
\newcommand\cvenfoque[1]{\def\cvEnfoque{#1}}
\newcommand\cvresumen[1]{\def\cvResumen{#1}}
\newcommand\cvubicacion[1]{\def\cvUbicacion{#1}}
\newcommand\cvemail[1]{\def\cvEmail{#1}}
\newcommand\cvlinkedin[1]{\def\cvLinkedin{#1}}
\newcommand\cvgithub[1]{\def\cvGithub{#1}}
\newcommand\cvorcid[1]{\def\cvOrcid{#1}}
\newcommand\cvweb[1]{\def\cvWeb{#1}}
\newcommand\cvtelefono[1]{\def\cvTelefono{#1}}
\newcommand\cvnacimiento[1]{\def\cvNacimiento{#1}}

% Section titles: \cvtitulos{clave=Título, ...} and \cvt{clave}.
\pgfkeys{/cvt/.unknown/.code={\csdef{cvt@\pgfkeyscurrentname}{#1}}}
\newcommand\cvtitulos[1]{\pgfqkeys{/cvt}{#1}}
\newcommand\cvt[1]{\csuse{cvt@#1}}

% --- Private data ----------------------------------------------------------
% privado.tex (git-ignored) holds the phone number and birth date. The public
% build (\cvpublico) leaves them out even when the file exists.
\providetoggle{cvpublico}
\ifdefined\cvpublic\toggletrue{cvpublico}\fi
\IfFileExists{privado.tex}{\input{privado.tex}}{}
\iftoggle{cvpublico}{\let\cvTelefono\relax\let\cvNacimiento\relax}{}
\newcommand\siTelefono[1]{\ifdefined\cvTelefono\ifx\cvTelefono\relax\else#1\fi\fi}
\newcommand\siNacimiento[1]{\ifdefined\cvNacimiento\ifx\cvNacimiento\relax\else#1\fi\fi}

% --- Flags -----------------------------------------------------------------
\newcommand\bandera[1]{%
  \begin{tikzpicture}[x=1mm,y=1mm,baseline=0.6mm]
    \ifstrequal{#1}{es}{%
      \fill[red!80!black] (0,0) rectangle (9,6);
      \fill[yellow!85!orange] (0,1.5) rectangle (9,4.5);
      \fill[red!70!black] (2,2.4) rectangle (3.2,3.6);
    }{%
      \fill[blue!50!black] (0,0) rectangle (9,6);
      \begin{scope}
        \clip (0,0) rectangle (9,6);
        \draw[white,line width=1.2mm] (0,0) -- (9,6) (0,6) -- (9,0);
        \draw[red!80!black,line width=0.4mm] (0,0) -- (9,6) (0,6) -- (9,0);
      \end{scope}
      \fill[white] (3.6,0) rectangle (5.4,6) (0,2.1) rectangle (9,3.9);
      \fill[red!80!black] (3.95,0) rectangle (5.05,6) (0,2.45) rectangle (9,3.55);
    }
  \end{tikzpicture}}

\setlist{nosep}

\usepackage{fontspec}
\usepackage{xcolor}

\providecommand\cvlang{es}

\defaultfontfeatures{Ligatures=TeX}
\setmainfont{Lora}[
  Path=../assets/fonts/, Extension=.ttf,
  UprightFont=*-Regular, BoldFont=*-Bold,
  ItalicFont=*-Italic, BoldItalicFont=*-BoldItalic]
\newfontfamily\lato{Lato}[
  Extension=.ttf, UprightFont=*-Regular, BoldFont=*-Bold]
\newfontfamily\latoheavy{Lato-Heavy.ttf}
\newfontfamily\latoblack{Lato-Black.ttf}

\definecolor{navy}{HTML}{2C3E50}
\definecolor{lateral}{HTML}{D9D9D9}
\definecolor{chip}{HTML}{5B7D95}
\definecolor{tinta}{HTML}{2B2B2B}

\pagestyle{empty}
\setlength\parindent{0pt}
\hyphenpenalty=10000\exhyphenpenalty=10000
\color{tinta}

\input{contenido-\cvlang}

% --- Main column pieces ----------------------------------------------------
\newcommand\cvbarra[1]{%
  \par\vspace{3.2mm}%
  \noindent\hspace*{-4.7mm}%
  {\setlength\fboxsep{0pt}\colorbox{navy}{\parbox[c][9mm]{140mm}{%
    \hspace{8.5mm}{\color{white}\latoblack\fontsize{12.5}{14}\selectfont
    \raisebox{0.25ex}{\tikz\fill[white] circle (1.05mm);}\hspace{2mm}%
    \addfontfeatures{LetterSpace=10}\MakeUppercase{#1}}}}}%
  \par\vspace{2.6mm}}

\newcommand\cvcentro[1]{%
  \par\vspace{1.2mm}\hspace*{8mm}%
  {\bfseries\fontsize{11.6}{14}\selectfont\addfontfeatures{LetterSpace=5}\MakeUppercase{#1}}\par\vspace{0.6mm}}

\newenvironment{cvlista}
  {\begin{itemize}[leftmargin=13.5mm,labelsep=2.2mm,itemsep=0.45mm,
      label={\normalfont\textbullet}]\raggedright\itshape\fontsize{10.6}{14.6}\selectfont\addfontfeatures{LetterSpace=2.5}}
  {\end{itemize}}

\newcommand\cvestudio[2]{#1 (#2)}

\newcommand\cvpuestotrabajo[3]{%
  \cvcentro{#1}%
  \hspace*{8mm}{\bfseries\itshape\fontsize{10.6}{14.6}\selectfont\addfontfeatures{LetterSpace=1.5}#2}\par
  {\raggedright\leftskip=8mm\itshape\fontsize{10.6}{14.6}\selectfont\addfontfeatures{LetterSpace=2.5}#3\par}%
  \vspace{2.4mm}}

\newcommand\cvdestacado[2]{{\upshape\bfseries #1} #2}
\newcommand\cvrol[2]{{\upshape\bfseries #1} {\upshape #2}}

% --- Sidebar pieces --------------------------------------------------------
\newcommand\cvidioma[3]{}
\newcommand\cvchip[1]{%
  \tikz[baseline=-0.6ex]\node[ellipse,fill=chip,text=white,inner sep=0.3mm,
    minimum height=7.4mm,minimum width=10.4mm,font=\latoblack,
    align=center]{\cvchiplabel{#1}};}
\newcommand\cvchiplabel[1]{%
  \StrLen{#1}[\cvlen]%
  \ifnum\cvlen<4 \fontsize{7.6}{8}\selectfont\MakeUppercase{#1}%
  \else\fontsize{4.4}{4.8}\selectfont
    \ifnum\cvlen>9 \parbox{8.6mm}{\centering\MakeUppercase{#1}}\else\MakeUppercase{#1}\fi\fi}
\usepackage{xstring}
\usetikzlibrary{shapes.geometric}
\newcommand\cvgrupotec[2]{%
  {\latoblack\fontsize{7.3}{9}\selectfont\addfontfeatures{LetterSpace=6}\MakeUppercase{#1}}\par
  \vspace{1.4mm}%
  {\baselineskip=0pt\lineskip=1.5mm\lineskiplimit=0pt\raggedright
   \forcsvlist{\cvchipsep}{#2}\par}\vspace{3.4mm}}
\newcommand\cvchipsep[1]{\cvchip{#1}\hspace{1.3mm}\allowbreak}

\newcommand\fila[3]{% y (mm from top), icon, text
  \node[anchor=base west,inner sep=0] at (8mm,-#1mm) {\color{black}\fontsize{9}{10}\selectfont#2};
  \node[anchor=base west,inner sep=0] at (12.6mm,-#1mm) {\lato\fontsize{8.6}{10}\selectfont\color{black!85}#3};}
\newcommand\enlace[2]{\href{#1}{\underline{#2}}}
\newcommand\titulolateral[2]{%
  \node[anchor=base west,inner sep=0] at (9mm,-#1mm)
    {\latoblack\fontsize{13}{15}\selectfont\addfontfeatures{LetterSpace=10}\color{black!90}\MakeUppercase{#2}};}
\newcommand\separador[1]{\draw[black,line width=0.55pt] (0,-#1mm) -- (72.9mm,-#1mm);}

\begin{document}
\vspace*{7.4mm}\noindent
\begin{tikzpicture}[remember picture,overlay,x=1mm,y=1mm]
  \begin{scope}[shift={(current page.north west)}]
    % Backgrounds.
    \fill[lateral] (0,0) rectangle (72.9,-297.1);
    \fill[navy] (0,0) rectangle (72.9,-77.2);
    \fill[navy] (72.9,0) rectangle (210.1,-52.2);

    % Portrait.
    \begin{scope}
      \clip (36.3,-27.2) circle (21.4);
      \node[inner sep=0] at (36.3,-29.2) {\includegraphics[width=43.6mm]{foto.jpg}};
    \end{scope}
    \draw[white!70!navy,line width=0.7pt] (36.3,-27.2) circle (21.4);

    % Name.
    \node[anchor=base,inner sep=0] at (36.3,-57.3)
      {\latoheavy\fontsize{19.5}{22}\selectfont\addfontfeatures{LetterSpace=4}\color{white}\MakeUppercase{\cvNombre}};
    \node[anchor=base,inner sep=0] at (36.3,-63.4)
      {\latoheavy\fontsize{17}{20}\selectfont\addfontfeatures{LetterSpace=4}\color{white}\MakeUppercase{\cvApellidos}};
    \node[anchor=base,inner sep=0] at (36.3,-71.7)
      {\lato\fontsize{12.4}{14}\selectfont\addfontfeatures{LetterSpace=2}\color{white}\cvPuesto};

    % Personal details.
    \fila{85.7}{\faMapMarker*}{\cvUbicacion}
    \siNacimiento{\fila{95.2}{\faCalendar[regular]}{\cvNacimiento}}
    \separador{101.0}
    \titulolateral{108.0}{\cvt{contacto}}
    \siTelefono{\fila{115.0}{\faPhone*}{\cvTelefono}}
    \fila{124.6}{\faEnvelope[regular]}{\enlace{mailto:\cvEmail}{\cvEmail}}
    \fila{134.2}{\faLinkedin}{\enlace{https://www.linkedin.com/in/\cvLinkedin}{\cvLinkedin}}
    \fila{143.8}{\faGithub}{\enlace{https://github.com/\cvGithub}{\cvGithub}}
    \ifdefined\cvOrcid\fila{153.4}{\aiOrcid}{\enlace{https://orcid.org/\cvOrcid}{\cvOrcid}}\fi
    \separador{160.7}
    \titulolateral{168.6}{\cvt{idiomas}}
    \renewcommand\cvidioma[3]{%
      \node[anchor=west,inner sep=0] at (\cvx,-178.2) {\bandera{#1}};
      \node[anchor=west,inner sep=0,align=left] at (\cvx+10.3,-178.2)
        {\lato\fontsize{8.6}{9.6}\selectfont#2\\#3};
      \pgfmathsetmacro\cvx{\cvx+24.6}\xdef\cvx{\cvx}}
    \def\cvx{6.8}
    \cvIdiomas
    \separador{185.9}
    \titulolateral{195.4}{\cvt{tecnologias}}
    \node[anchor=north west,inner sep=0,text width=66mm] at (5.2,-198.8) {\cvTecnologias};
  \end{scope}
\end{tikzpicture}%
{\color{white}\tolerance=9999\emergencystretch=4em\fontsize{12.4}{18.6}\selectfont\addfontfeatures{LetterSpace=2.5}\cvResumen\par}
\vspace{-1.6mm}
\cvbarra{\cvt{formacion}}
\cvFormacion
\cvbarra{\cvt{aptitudes}}
\cvAptitudes
\cvbarra{\cvt{experiencia}}
\cvExperiencia
\vspace{-2.4mm}
\cvbarra{\cvt{proyectos}}
\cvProyectos
\cvbarra{\cvt{complementaria}}
\cvComplementaria
\end{document}
"   (o: node build.mjs)
\documentclass[10pt]{article}
\usepackage[a4paper,left=77.7mm,right=7mm,top=0mm,bottom=0mm]{geometry}
% Shared by both templates: storage for the content fields, the private data,
% flags and icons. Compile with XeLaTeX (see build.mjs).
\usepackage{etoolbox}
\usepackage{pgfkeys}
\usepackage{graphicx}
\usepackage{tikz}
\usepackage{fontawesome5}
\usepackage{academicons}
\usepackage{enumitem}
\usepackage[hidelinks]{hyperref}

% --- Content fields --------------------------------------------------------
\newcommand\cvnombre[2]{\def\cvNombre{#1}\def\cvApellidos{#2}}
\newcommand\cvpuesto[1]{\def\cvPuesto{#1}}
\newcommand\cvenfoque[1]{\def\cvEnfoque{#1}}
\newcommand\cvresumen[1]{\def\cvResumen{#1}}
\newcommand\cvubicacion[1]{\def\cvUbicacion{#1}}
\newcommand\cvemail[1]{\def\cvEmail{#1}}
\newcommand\cvlinkedin[1]{\def\cvLinkedin{#1}}
\newcommand\cvgithub[1]{\def\cvGithub{#1}}
\newcommand\cvorcid[1]{\def\cvOrcid{#1}}
\newcommand\cvweb[1]{\def\cvWeb{#1}}
\newcommand\cvtelefono[1]{\def\cvTelefono{#1}}
\newcommand\cvnacimiento[1]{\def\cvNacimiento{#1}}

% Section titles: \cvtitulos{clave=Título, ...} and \cvt{clave}.
\pgfkeys{/cvt/.unknown/.code={\csdef{cvt@\pgfkeyscurrentname}{#1}}}
\newcommand\cvtitulos[1]{\pgfqkeys{/cvt}{#1}}
\newcommand\cvt[1]{\csuse{cvt@#1}}

% --- Private data ----------------------------------------------------------
% privado.tex (git-ignored) holds the phone number and birth date. The public
% build (\cvpublico) leaves them out even when the file exists.
\providetoggle{cvpublico}
\ifdefined\cvpublic\toggletrue{cvpublico}\fi
\IfFileExists{privado.tex}{\input{privado.tex}}{}
\iftoggle{cvpublico}{\let\cvTelefono\relax\let\cvNacimiento\relax}{}
\newcommand\siTelefono[1]{\ifdefined\cvTelefono\ifx\cvTelefono\relax\else#1\fi\fi}
\newcommand\siNacimiento[1]{\ifdefined\cvNacimiento\ifx\cvNacimiento\relax\else#1\fi\fi}

% --- Flags -----------------------------------------------------------------
\newcommand\bandera[1]{%
  \begin{tikzpicture}[x=1mm,y=1mm,baseline=0.6mm]
    \ifstrequal{#1}{es}{%
      \fill[red!80!black] (0,0) rectangle (9,6);
      \fill[yellow!85!orange] (0,1.5) rectangle (9,4.5);
      \fill[red!70!black] (2,2.4) rectangle (3.2,3.6);
    }{%
      \fill[blue!50!black] (0,0) rectangle (9,6);
      \begin{scope}
        \clip (0,0) rectangle (9,6);
        \draw[white,line width=1.2mm] (0,0) -- (9,6) (0,6) -- (9,0);
        \draw[red!80!black,line width=0.4mm] (0,0) -- (9,6) (0,6) -- (9,0);
      \end{scope}
      \fill[white] (3.6,0) rectangle (5.4,6) (0,2.1) rectangle (9,3.9);
      \fill[red!80!black] (3.95,0) rectangle (5.05,6) (0,2.45) rectangle (9,3.55);
    }
  \end{tikzpicture}}

\setlist{nosep}

\usepackage{fontspec}
\usepackage{xcolor}

\providecommand\cvlang{es}

\defaultfontfeatures{Ligatures=TeX}
\setmainfont{Lora}[
  Path=../assets/fonts/, Extension=.ttf,
  UprightFont=*-Regular, BoldFont=*-Bold,
  ItalicFont=*-Italic, BoldItalicFont=*-BoldItalic]
\newfontfamily\lato{Lato}[
  Extension=.ttf, UprightFont=*-Regular, BoldFont=*-Bold]
\newfontfamily\latoheavy{Lato-Heavy.ttf}
\newfontfamily\latoblack{Lato-Black.ttf}

\definecolor{navy}{HTML}{2C3E50}
\definecolor{lateral}{HTML}{D9D9D9}
\definecolor{chip}{HTML}{5B7D95}
\definecolor{tinta}{HTML}{2B2B2B}

\pagestyle{empty}
\setlength\parindent{0pt}
\hyphenpenalty=10000\exhyphenpenalty=10000
\color{tinta}

\input{contenido-\cvlang}

% --- Main column pieces ----------------------------------------------------
\newcommand\cvbarra[1]{%
  \par\vspace{3.2mm}%
  \noindent\hspace*{-4.7mm}%
  {\setlength\fboxsep{0pt}\colorbox{navy}{\parbox[c][9mm]{140mm}{%
    \hspace{8.5mm}{\color{white}\latoblack\fontsize{12.5}{14}\selectfont
    \raisebox{0.25ex}{\tikz\fill[white] circle (1.05mm);}\hspace{2mm}%
    \addfontfeatures{LetterSpace=10}\MakeUppercase{#1}}}}}%
  \par\vspace{2.6mm}}

\newcommand\cvcentro[1]{%
  \par\vspace{1.2mm}\hspace*{8mm}%
  {\bfseries\fontsize{11.6}{14}\selectfont\addfontfeatures{LetterSpace=5}\MakeUppercase{#1}}\par\vspace{0.6mm}}

\newenvironment{cvlista}
  {\begin{itemize}[leftmargin=13.5mm,labelsep=2.2mm,itemsep=0.45mm,
      label={\normalfont\textbullet}]\raggedright\itshape\fontsize{10.6}{14.6}\selectfont\addfontfeatures{LetterSpace=2.5}}
  {\end{itemize}}

\newcommand\cvestudio[2]{#1 (#2)}

\newcommand\cvpuestotrabajo[3]{%
  \cvcentro{#1}%
  \hspace*{8mm}{\bfseries\itshape\fontsize{10.6}{14.6}\selectfont\addfontfeatures{LetterSpace=1.5}#2}\par
  {\raggedright\leftskip=8mm\itshape\fontsize{10.6}{14.6}\selectfont\addfontfeatures{LetterSpace=2.5}#3\par}%
  \vspace{2.4mm}}

\newcommand\cvdestacado[2]{{\upshape\bfseries #1} #2}
\newcommand\cvrol[2]{{\upshape\bfseries #1} {\upshape #2}}

% --- Sidebar pieces --------------------------------------------------------
\newcommand\cvidioma[3]{}
\newcommand\cvchip[1]{%
  \tikz[baseline=-0.6ex]\node[ellipse,fill=chip,text=white,inner sep=0.3mm,
    minimum height=7.4mm,minimum width=10.4mm,font=\latoblack,
    align=center]{\cvchiplabel{#1}};}
\newcommand\cvchiplabel[1]{%
  \StrLen{#1}[\cvlen]%
  \ifnum\cvlen<4 \fontsize{7.6}{8}\selectfont\MakeUppercase{#1}%
  \else\fontsize{4.4}{4.8}\selectfont
    \ifnum\cvlen>9 \parbox{8.6mm}{\centering\MakeUppercase{#1}}\else\MakeUppercase{#1}\fi\fi}
\usepackage{xstring}
\usetikzlibrary{shapes.geometric}
\newcommand\cvgrupotec[2]{%
  {\latoblack\fontsize{7.3}{9}\selectfont\addfontfeatures{LetterSpace=6}\MakeUppercase{#1}}\par
  \vspace{1.4mm}%
  {\baselineskip=0pt\lineskip=1.5mm\lineskiplimit=0pt\raggedright
   \forcsvlist{\cvchipsep}{#2}\par}\vspace{3.4mm}}
\newcommand\cvchipsep[1]{\cvchip{#1}\hspace{1.3mm}\allowbreak}

\newcommand\fila[3]{% y (mm from top), icon, text
  \node[anchor=base west,inner sep=0] at (8mm,-#1mm) {\color{black}\fontsize{9}{10}\selectfont#2};
  \node[anchor=base west,inner sep=0] at (12.6mm,-#1mm) {\lato\fontsize{8.6}{10}\selectfont\color{black!85}#3};}
\newcommand\enlace[2]{\href{#1}{\underline{#2}}}
\newcommand\titulolateral[2]{%
  \node[anchor=base west,inner sep=0] at (9mm,-#1mm)
    {\latoblack\fontsize{13}{15}\selectfont\addfontfeatures{LetterSpace=10}\color{black!90}\MakeUppercase{#2}};}
\newcommand\separador[1]{\draw[black,line width=0.55pt] (0,-#1mm) -- (72.9mm,-#1mm);}

\begin{document}
\vspace*{7.4mm}\noindent
\begin{tikzpicture}[remember picture,overlay,x=1mm,y=1mm]
  \begin{scope}[shift={(current page.north west)}]
    % Backgrounds.
    \fill[lateral] (0,0) rectangle (72.9,-297.1);
    \fill[navy] (0,0) rectangle (72.9,-77.2);
    \fill[navy] (72.9,0) rectangle (210.1,-52.2);

    % Portrait.
    \begin{scope}
      \clip (36.3,-27.2) circle (21.4);
      \node[inner sep=0] at (36.3,-29.2) {\includegraphics[width=43.6mm]{foto.jpg}};
    \end{scope}
    \draw[white!70!navy,line width=0.7pt] (36.3,-27.2) circle (21.4);

    % Name.
    \node[anchor=base,inner sep=0] at (36.3,-57.3)
      {\latoheavy\fontsize{19.5}{22}\selectfont\addfontfeatures{LetterSpace=4}\color{white}\MakeUppercase{\cvNombre}};
    \node[anchor=base,inner sep=0] at (36.3,-63.4)
      {\latoheavy\fontsize{17}{20}\selectfont\addfontfeatures{LetterSpace=4}\color{white}\MakeUppercase{\cvApellidos}};
    \node[anchor=base,inner sep=0] at (36.3,-71.7)
      {\lato\fontsize{12.4}{14}\selectfont\addfontfeatures{LetterSpace=2}\color{white}\cvPuesto};

    % Personal details.
    \fila{85.7}{\faMapMarker*}{\cvUbicacion}
    \siNacimiento{\fila{95.2}{\faCalendar[regular]}{\cvNacimiento}}
    \separador{101.0}
    \titulolateral{108.0}{\cvt{contacto}}
    \siTelefono{\fila{115.0}{\faPhone*}{\cvTelefono}}
    \fila{124.6}{\faEnvelope[regular]}{\enlace{mailto:\cvEmail}{\cvEmail}}
    \fila{134.2}{\faLinkedin}{\enlace{https://www.linkedin.com/in/\cvLinkedin}{\cvLinkedin}}
    \fila{143.8}{\faGithub}{\enlace{https://github.com/\cvGithub}{\cvGithub}}
    \ifdefined\cvOrcid\fila{153.4}{\aiOrcid}{\enlace{https://orcid.org/\cvOrcid}{\cvOrcid}}\fi
    \separador{160.7}
    \titulolateral{168.6}{\cvt{idiomas}}
    \renewcommand\cvidioma[3]{%
      \node[anchor=west,inner sep=0] at (\cvx,-178.2) {\bandera{#1}};
      \node[anchor=west,inner sep=0,align=left] at (\cvx+10.3,-178.2)
        {\lato\fontsize{8.6}{9.6}\selectfont#2\\#3};
      \pgfmathsetmacro\cvx{\cvx+24.6}\xdef\cvx{\cvx}}
    \def\cvx{6.8}
    \cvIdiomas
    \separador{185.9}
    \titulolateral{195.4}{\cvt{tecnologias}}
    \node[anchor=north west,inner sep=0,text width=66mm] at (5.2,-198.8) {\cvTecnologias};
  \end{scope}
\end{tikzpicture}%
{\color{white}\tolerance=9999\emergencystretch=4em\fontsize{12.4}{18.6}\selectfont\addfontfeatures{LetterSpace=2.5}\cvResumen\par}
\vspace{-1.6mm}
\cvbarra{\cvt{formacion}}
\cvFormacion
\cvbarra{\cvt{aptitudes}}
\cvAptitudes
\cvbarra{\cvt{experiencia}}
\cvExperiencia
\vspace{-2.4mm}
\cvbarra{\cvt{proyectos}}
\cvProyectos
\cvbarra{\cvt{complementaria}}
\cvComplementaria
\end{document}
